%======================================================================
% InsertSplit: the one-site quotient of the blocked-region weight
%======================================================================

\begin{theorem}[Region vertex product across an inserted site]%
    \label{thm:peps_prod_region_insert_split}
    \lean{TNLean.PEPS.prod_region_insert_split}
    \leanok
    Let $v\notin R$.  At a fixed global virtual configuration $\zeta$ and a
    physical configuration $\sigma$ on $R\cup\{v\}$, the product of the local
    tensors over the vertices of $R\cup\{v\}$ factors as the local tensor at the
    inserted site $v$ times the product over the vertices of $R$:
    \[
        \prod_{w\in R\cup\{v\}}A_w\bigl(\zeta|_{\operatorname{inc}(w)},\sigma_w\bigr)
        =
        A_v\bigl(\zeta|_{\operatorname{inc}(v)},\sigma_v\bigr)
        \prod_{w\in R}A_w\bigl(\zeta|_{\operatorname{inc}(w)},\sigma|_R{}_w\bigr).
    \]
\end{theorem}
\begin{proof}\leanok
    The inserted site $v$ is isolated from the product over $R\cup\{v\}$ by
    splitting off a single factor, and the residual product over the remaining
    vertices is reindexed through the bijection between the vertices of $R$ and
    the non-$v$ vertices of $R\cup\{v\}$.
\end{proof}

\begin{theorem}[One-site split of the blocked-region weight]%
    \label{thm:peps_regionBlockedWeight_insert_eq_sum_split}
    \lean{TNLean.PEPS.regionBlockedWeight_insert_eq_sum_split}
    \leanok
    \uses{thm:peps_prod_region_insert_split}
    Let $v\notin R$.  The blocked-region weight of $R\cup\{v\}$ is the constrained
    sum, over global virtual configurations $\zeta$ whose crossing-edge label on
    $R\cup\{v\}$ is $\mu$, of the inserted-site tensor times the vertex product
    over $R$:
    \[
        T_{A,R\cup\{v\}}^{\mu}(\sigma)
        =
        \sum_{\zeta:\,\zeta|_{\partial(R\cup\{v\})}=\mu}
        A_v\bigl(\zeta|_{\operatorname{inc}(v)},\sigma_v\bigr)
        \prod_{w\in R}A_w\bigl(\zeta|_{\operatorname{inc}(w)},\sigma|_R{}_w\bigr).
    \]
\end{theorem}
\begin{proof}\leanok
    Each summand of the blocked-region weight is the vertex product over
    $R\cup\{v\}$, which factors by
    Theorem~\ref{thm:peps_prod_region_insert_split}.
\end{proof}

\begin{theorem}[A non-incident crossing edge survives the insertion]%
    \label{thm:peps_isRegionBoundaryEdge_insert_of_not_incident}
    \lean{TNLean.PEPS.isRegionBoundaryEdge_insert_of_not_incident}
    \leanok
    Let $v\notin R$.  A crossing edge of $R$ whose two endpoints both differ from
    $v$ is also a crossing edge of $R\cup\{v\}$.
\end{theorem}
\begin{proof}\leanok
    A crossing edge of $R$ has exactly one endpoint inside $R$.  Adjoining $v$ to
    $R$ leaves both endpoints on the same side as before, since neither endpoint
    is $v$, so the edge still crosses the boundary of $R\cup\{v\}$.
\end{proof}

\begin{definition}[Boundary label of the smaller region across an insertion]%
    \label{def:peps_boundaryLabelOfInsert}
    \lean{TNLean.PEPS.boundaryLabelOfInsert}
    \leanok
    \uses{thm:peps_isRegionBoundaryEdge_insert_of_not_incident}
    Let $v\notin R$.  Given a crossing-edge label $\mu$ for $R\cup\{v\}$ and a
    local virtual configuration $\eta$ on the edges incident to $v$, the
    crossing-edge label $\beta(\mu,\eta)$ of $R$ is read off these two
    configurations: a crossing edge of $R$ incident to $v$ has become internal to
    $R\cup\{v\}$, so its label is taken from $\eta$; a crossing edge of $R$ not
    incident to $v$ is also a crossing edge of $R\cup\{v\}$, so its label is taken
    from $\mu$.
\end{definition}

\begin{theorem}[Boundary-label fiber identity across an inserted site]%
    \label{thm:peps_regionBoundaryLabel_eq_boundaryLabelOfInsert}
    \lean{TNLean.PEPS.regionBoundaryLabel_eq_boundaryLabelOfInsert}
    \leanok
    \uses{def:peps_boundaryLabelOfInsert}
    Let $v\notin R$.  If a global virtual configuration $\zeta$ restricts to $\mu$
    on the crossing edges of $R\cup\{v\}$ and to $\eta$ on the edges incident to
    $v$, then the crossing-edge label of $R$ read from $\zeta$ is exactly
    $\beta(\mu,\eta)$ of
    Definition~\ref{def:peps_boundaryLabelOfInsert}.
\end{theorem}
\begin{proof}\leanok
    The two labels are compared edge by edge.  On a crossing edge of $R$
    incident to $v$ both sides read $\zeta$ on an edge incident to $v$, which is
    $\eta$.  On a crossing edge of $R$ not incident to $v$ both sides read
    $\zeta$ on a crossing edge of $R\cup\{v\}$, which is $\mu$.
\end{proof}

%======================================================================
% Gauge absorption at region granularity
%======================================================================

\begin{theorem}[Region gauge absorption]%
    \label{thm:peps_regionInsertedCoeff_applyGauge}
    \lean{TNLean.PEPS.regionInsertedCoeff_applyGauge}
    \leanok
    \uses{thm:peps_regionInsertedCoeff_eq_smul_edgeInsertedCoeff,
        thm:peps_edge_gauge_absorption, def:applyGauge}
    Let $X$ be an edge-gauge family, write $B^X$ for the gauged tensor, and let
    $f$ be a crossing edge of $R$.  Inserting $M$ on $f$ between the region and
    its complement in $B^X$ equals inserting, in $B$, the matrix obtained by
    orienting $M$ to the boundary-edge convention, conjugating by the transpose
    of the endpoint gauge $X_f$, and orienting back:
    \[
        C_{B^X}(R,f,M;\sigma,\tau)
        =
        C_B\bigl(R,f,
            \mathcal O_{R,f}\bigl(X_f^{\mathsf T}\,
                \mathcal O_{R,f}(M)\,(X_f^{-1})^{\mathsf T}\bigr);
            \sigma,\tau\bigr),
    \]
    where $\mathcal O_{R,f}$ is the boundary-edge orientation.  This is the
    region-granularity analogue of
    Theorem~\ref{thm:peps_edge_gauge_absorption}.
\end{theorem}
\begin{proof}\leanok
    By Theorem~\ref{thm:peps_regionInsertedCoeff_eq_smul_edgeInsertedCoeff} the
    region/complement contraction is the single-bond cut at $f$ up to the
    interior-bond multiplicity $p_B(R)=\prod_{g\notin\partial R}D_B(g)$, which is
    unchanged by a gauge.  At the edge granularity every interior edge and every
    crossing edge other than $f$ cancels its gauge pairwise, while on $f$ the
    endpoint gauges conjugate the oriented matrix
    $N=\mathcal O_{R,f}(M)$ by the transpose of $X_f$ and its inverse:
    \[
        E_{B^X}\bigl(f,\omega_{R,\sigma,\tau},N\bigr)
        =
        E_B\bigl(f,\omega_{R,\sigma,\tau},
            X_f^{\mathsf T}\,N\,(X_f^{-1})^{\mathsf T}\bigr).
    \]
    Transporting this back through
    Theorem~\ref{thm:peps_regionInsertedCoeff_eq_smul_edgeInsertedCoeff} on $B$,
    and using that the boundary-edge orientation is an involution
    ($\mathcal O_{R,f}\circ\mathcal O_{R,f}=\operatorname{id}$), gives the
    displayed equality.
\end{proof}

\subsection{From the conjugation identity to the absorbed equality}%
\label{subsec:peps_absorbed_equality}

% Companion to docs/paper-gaps/peps_normal_ft_section3_route.tex.
At a region $R$ with boundary edge $f$, the per-edge gauge identity takes
the \emph{conjugation} form
\[
    C_A(R,f,M;\sigma,\tau)
    =
    C_B\bigl(R,f,Z\,\phi_{h_f}(M)\,Z^{-1};\sigma,\tau\bigr),
\]
where $Z$ is the invertible bond matrix on $f$ and $\phi_{h_f}$ is the
matrix-algebra transport induced by $h_f:D_A(f)=D_B(f)$. The region comparison
of \cite[Section~3]{Molnar2018NormalPEPS} instead consumes the \emph{absorbed}
form: the gauge is absorbed into the second tensor, and the \emph{same} matrix
$\phi_{h_f}(M)$ is inserted on both sides.

\begin{definition}[Orientation-adapted absorbing gauge]%
    \label{def:peps_absorbedBoundaryGauge}
    \lean{TNLean.PEPS.absorbedBoundaryGauge}
    \leanok
    Let $f$ be a boundary edge of $R$ and let $Z$ be an invertible matrix on the
    bond space of $f$. The orientation-adapted absorbing gauge is
    \[
        X_f
        =
        \begin{cases}
            Z^{\mathsf T} & \text{if the left endpoint of $f$ lies in $R$,}\\
            Z^{-1} & \text{otherwise.}
        \end{cases}
    \]
    It is the absorbing matrix on $f$ for which the gauge action on the region
    coefficient reproduces the conjugation $Z\,(\cdot)\,Z^{-1}$.
\end{definition}

\begin{theorem}[Matrix reconciliation for the absorbing gauge]%
    \label{thm:peps_regionEdgeOrient_absorbedBoundaryGauge}
    \lean{TNLean.PEPS.regionEdgeOrient_absorbedBoundaryGauge}
    \leanok
    \uses{def:peps_absorbedBoundaryGauge}
    With $X_f$ the orientation-adapted absorbing gauge of $Z$ and
    $\mathcal O_{R,f}$ the boundary-edge orientation, for every matrix $N$ on the
    bond space of $f$,
    \[
        \mathcal O_{R,f}\bigl(
            X_f^{\mathsf T}\,\mathcal O_{R,f}(N)\,(X_f^{-1})^{\mathsf T}
        \bigr)
        =
        Z\,N\,Z^{-1}.
    \]
    The conjugation by $X_f^{\mathsf T}\cdot(\cdot)\cdot(X_f^{-1})^{\mathsf T}$
    carried by the region gauge action thus equals $Z\,(\cdot)\,Z^{-1}$.
\end{theorem}
\begin{proof}\leanok
    When the left endpoint of $f$ lies in $R$, the orientation is the identity
    and $X_f=Z^{\mathsf T}$, so $X_f^{\mathsf T}=Z$ and
    $(X_f^{-1})^{\mathsf T}=Z^{-1}$. Otherwise the orientation is transpose and
    $X_f=Z^{-1}$, so $X_f^{\mathsf T}=(Z^{-1})^{\mathsf T}$ and
    $(X_f^{-1})^{\mathsf T}=Z^{\mathsf T}$, giving
    \[
        X_f^{\mathsf T}\,\mathcal O_{R,f}(N)\,(X_f^{-1})^{\mathsf T}
        =
        (Z^{-1})^{\mathsf T}\,N^{\mathsf T}\,Z^{\mathsf T}
        =
        (Z\,N\,Z^{-1})^{\mathsf T}.
    \]
    Applying the outer transpose orientation gives
    $\mathcal O_{R,f}((Z\,N\,Z^{-1})^{\mathsf T})=Z\,N\,Z^{-1}$.
    Both cases give $Z\,N\,Z^{-1}$.
\end{proof}

\begin{theorem}[Region absorbed equality from the conjugation identity]%
    \label{thm:peps_regionInsertedCoeff_eq_applyGauge_of_conjIdentity}
    \lean{TNLean.PEPS.regionInsertedCoeff_eq_applyGauge_of_conjIdentity}
    \leanok
    \uses{def:peps_regionInsertedCoeff, def:applyGauge,
        def:peps_absorbedBoundaryGauge,
        thm:peps_regionEdgeOrient_absorbedBoundaryGauge,
        thm:peps_regionInsertedCoeff_applyGauge}
    Let $R$ be a region with boundary edge $f$ and assume the
    bond-dimension equality $h:D_A=D_B$. Suppose the per-edge gauge $Z$ realizes
    the conjugation-form region-insertion identity at $f$: for every matrix $M$
    on the bond space of $f$ and all physical configurations $\sigma$ on $R$ and
    $\tau$ on $V\setminus R$,
    \[
        C_A(R,f,M;\sigma,\tau)
        =
        C_B\bigl(R,f,Z\,\phi_{h_f}(M)\,Z^{-1};\sigma,\tau\bigr).
    \]
    Let $X$ be an edge-gauge family whose value on $f$ is the orientation-adapted
    absorbing gauge $X_f$ of $Z$. Then $A$ inserting $M$ matches the absorbed
    tensor $B^X$ inserting the same transported matrix:
    \[
        C_A(R,f,M;\sigma,\tau)
        =
        C_{B^X}\bigl(R,f,\phi_{h_f}(M);\sigma,\tau\bigr).
    \]
    This is the region analogue of
    Theorem~\ref{thm:peps_postAbsorptionEdgeInsertionEquality}: the
    conjugation-form coefficient identity becomes the absorbed plain equality.
\end{theorem}
\begin{proof}\leanok
    By Theorem~\ref{thm:peps_regionInsertedCoeff_applyGauge} the region gauge
    action conjugates the inserted matrix by
    $X_f^{\mathsf T}\cdot\mathcal O_{R,f}(\cdot)\cdot(X_f^{-1})^{\mathsf T}$ and
    cancels every interior gauge of the region and of its complement, so the
    right-hand coefficient over $B^X$ equals the coefficient over $B$ with the
    transported matrix conjugated by that gauge. Substituting the conjugation
    hypothesis on the left and the matrix reconciliation
    Theorem~\ref{thm:peps_regionEdgeOrient_absorbedBoundaryGauge} matches the two
    conjugation forms.
\end{proof}

\begin{theorem}[Edge absorbed equality from the region absorbed equality]%
    \label{thm:peps_edgeInsertedCoeff_eq_applyGauge_of_region}
    \lean{TNLean.PEPS.edgeInsertedCoeff_eq_applyGauge_of_region}
    \leanok
    \uses{def:peps_regionInsertedCoeff, def:peps_edgeInsertedCoeff,
        def:applyGauge,
        thm:peps_regionInsertedCoeff_eq_smul_edgeInsertedCoeff}
    Let $R$ be a region with boundary edge $f$, assume the bond-dimension
    equality $h:D_A=D_B$ and that every bond dimension of $A$ is positive, and
    suppose the absorbed region equality holds: for every matrix $M$ on the bond
    space of $f$ and all physical configurations $\sigma$ on $R$ and $\tau$ on
    $V\setminus R$,
    \[
        C_A(R,f,M;\sigma,\tau)
        =
        C_{B^X}\bigl(R,f,\phi_{h_f}(M);\sigma,\tau\bigr).
    \]
    Then the edge-level absorbed equality holds at $f$: for every global physical
    configuration $\sigma$ and every matrix $N$ on the bond space of $f$,
    \[
        E_A(f,\sigma,N)
        =
        E_{B^X}\bigl(f,\sigma,\phi_{h_f}(N)\bigr).
    \]
    The edge-level identity no longer mentions the region.
\end{theorem}
\begin{proof}\leanok
    The interior bond multiplicity $p_A(R)=\prod_{g\notin\partial R}D_A(g)$ is a
    positive integer, and it is shared by $A$ and $B^X$ because the gauge
    preserves bond dimensions and $D_A=D_B$. Apply
    Theorem~\ref{thm:peps_regionInsertedCoeff_eq_smul_edgeInsertedCoeff} with the
    inserted matrix $\mathcal O_{R,f}(N)$ and with configurations
    $(\sigma|_R,\sigma|_{V\setminus R})$. On the left-hand side,
    \[
        C_A\bigl(R,f,\mathcal O_{R,f}(N);
            \sigma|_R,\sigma|_{V\setminus R}\bigr)
        =
        p_A(R)\,
        E_A\bigl(f,\sigma,
            \mathcal O_{R,f}(\mathcal O_{R,f}(N))\bigr)
        =
        p_A(R)\,E_A(f,\sigma,N).
    \]
    On the right-hand side of the assumed region equality,
    \[
    \begin{aligned}
        C_{B^X}\bigl(R,f,\phi_{h_f}(\mathcal O_{R,f}(N));
            \sigma|_R,\sigma|_{V\setminus R}\bigr)
        &=
        p_A(R)\,
        E_{B^X}\bigl(f,\sigma,
            \mathcal O_{R,f}(\phi_{h_f}(\mathcal O_{R,f}(N)))\bigr)\\
        &=
        p_A(R)\,
        E_{B^X}\bigl(f,\sigma,\phi_{h_f}(N)\bigr).
    \end{aligned}
    \]
    The displayed equalities also use
    $\mathcal O_{R,f}\circ\mathcal O_{R,f}=\operatorname{id}$, the compatibility
    of $\mathcal O_{R,f}$ with transport, and
    $\omega_{R,\sigma|_R,\sigma|_{V\setminus R}}=\sigma$. Substituting these two
    identities into the assumed region equality and cancelling the shared
    positive multiplicity gives the edge-level identity.
\end{proof}

\begin{theorem}[Region absorbed equality from the edge absorbed equality]%
    \label{thm:peps_regionInsertedCoeff_eq_applyGauge_of_edge}
    \lean{TNLean.PEPS.regionInsertedCoeff_eq_applyGauge_of_edge}
    \leanok
    \uses{def:peps_regionInsertedCoeff, def:peps_edgeInsertedCoeff,
        def:applyGauge,
        thm:peps_regionInsertedCoeff_eq_smul_edgeInsertedCoeff}
    Assume the bond-dimension equality $h:D_A=D_B$ and suppose the edge-level
    absorbed equality holds at an edge $e$: for every global physical
    configuration $\sigma$ and every matrix $N$ on the bond space of $e$,
    \[
        E_A(e,\sigma,N)
        =
        E_{B^X}\bigl(e,\sigma,\phi_{h_e}(N)\bigr).
    \]
    Then the region absorbed equality holds at every region $R$ having $e$ as a
    boundary edge: for every matrix $M$ on the bond space of $e$ and all physical
    configurations $\sigma$ on $R$ and $\tau$ on $V\setminus R$,
    \[
        C_A(R,e,M;\sigma,\tau)
        =
        C_{B^X}\bigl(R,e,\phi_{h_e}(M);\sigma,\tau\bigr).
    \]
    Together with the preceding theorem this expresses region-independence: the
    absorbed plain equality over any comparison region follows from the single
    edge-level identity at the boundary edge.
\end{theorem}
\begin{proof}\leanok
    Read both region coefficients through the region-to-edge identity
    Theorem~\ref{thm:peps_regionInsertedCoeff_eq_smul_edgeInsertedCoeff}. The
    interior multiplicities agree because the gauge preserves bond dimensions and
    $D_A=D_B$, and the orientation $\mathcal O_{R,e}$ commutes with the transport.
    Hence the two sides are $p_A(R)$ times the edge coefficients
    \[
        E_A(e,\omega_{R,\sigma,\tau},\mathcal O_{R,e}(M))
        \quad\text{and}\quad
        E_{B^X}\bigl(e,\omega_{R,\sigma,\tau},
            \phi_{h_e}(\mathcal O_{R,e}(M))\bigr).
    \]
    The edge-level hypothesis at $\omega_{R,\sigma,\tau}$ equates these
    coefficients; multiplying by the common factor $p_A(R)$ gives the region
    equality.
\end{proof}

